\chapter{المقدمة}

\section{التمهيد والخلفية العلمية}
تلعب الأذرع الروبوتية دوراً محورياً في نهضة الأتمتة الصناعية الحديثة، حيث تُوكل إليها مهام معقدة تتطلب دقة متناهية وسرعة استجابة في بيئات العمل المختلفة. من منظور هندسة التحكم الآلي، تُعد الذراع الروبوتية نظاماً ميكانيكياً شديد اللاخطية (\textenglish{Highly Non-linear System})، تخضع حركتها لتأثيرات ديناميكية معقدة تشمل القصور الذاتي، وقوى الجاذبية، وتغيرات الحمولة. لذلك، فإن التحدي الهندسي الأبرز لا يقتصر فقط على بناء الهيكل الميكانيكي، بل يمتد إلى اشتقاق نماذج رياضية دقيقة (حركية وديناميكية) وتصميم خوارزميات تحكم متينة قادرة على ضمان استقرار النظام ودقة تتبعه للمسارات المطلوبة في الفضاء الكارتزي.

إلى جانب هذه التحديات في طبقات التحكم منخفض المستوى (\textenglish{LLC})، يتجه التطور الحديث في الروبوتيك نحو تحسين واجهات التفاعل البشري-الروبوتي (\textenglish{Human-Robot Interaction - HRI}). بدلاً من الاعتماد الحصري على لوحات التحكم التقليدية أو الواجهات البرمجية المعقدة، برزت تقنيات "التعرف على الكلام" (\textenglish{Speech Recognition}) كطبقة تحكم إشرافية عُليا تتيح توجيه الروبوت بمرونة وفعالية، مما يحرر يدي المشغل ويقلل من عبء التفاعل الميكانيكي، لتصبح المنظومة أكثر ذكاءً وقابلية للدمج في التطبيقات الصناعية والخدمية المعاصرة.

\section{دوافع المشروع وأهميته}
تنبثق أهمية هذا المشروع من الحاجة إلى بناء منظومة روبوتية متكاملة هندسياً، لا تكتفي بتنفيذ الحركة رياضيّاً فحسب، بل تضمن أداءً ديناميكياً مستقراً في مواجهة التحديات التشغيلية الواقعية. يتمثل الدافع الأساسي للعمل في التغلب على مشاكل التحكم الكلاسيكي عند التعامل مع الأنظمة الروبوتية متعددة المفاصل، وتحديداً التحدي المتمثل في التعامل مع التغيرات المفاجئة في الخصائص الفيزيائية للنظام، مثل تغير الكتلة وعزوم القصور الذاتي الناتجة عن التقاط ونقل الحمولات المختلفة أثناء التشغيل.

بناءً على ذلك، يسعى المشروع لتقديم حل رياضي وتنفيذي متين عبر تطوير وتقييم استراتيجيات تحكم متقدمة مثل التحكم التكيفي (\textenglish{ACTC}) والتحكم الانزلاقي (\textenglish{SMC}) المتغيرة والاضطرابات الديناميكية، مما يضمن دقة عالية في تتبع المسار الكارتزي ويعزز من كفاءة وموثوقية الأذرع الروبوتية في التطبيقات العملية.
\section{أهداف المشروع}
يهدف هذا المشروع إلى تصميم وتنفيذ ذراع روبوتية ذات خمس درجات حرية (\textenglish{5-DOF}) مدعومة بنظام تحكم هرمي متكامل. ولتحقيق ذلك، تم تحديد الأهداف المرحلية التالية:
\begin{enumerate}
    \item \textbf{النمذجة الرياضية الدقيقة:} اشتقاق النماذج الحركية (الأمامية والعكسية) باستخدام قواعد دينافيت-هارتنبرغ (\textenglish{DH})، وتأسيس النموذج الديناميكي وفق صياغة لاغرانج (\textenglish{Lagrange}) لحساب مصفوفات العطالة والجاذبية.
    \item \textbf{تصميم وتطوير خوارزميات التحكم:} تصميم واختبار عدة خوارزميات للتحكم منخفض المستوى (\textenglish{LLC})، مقارنة أداء متحكم \textenglish{PID} مع التحكم التكيفي (\textenglish{ACTC}) والتحكم الانزلاقي (\textenglish{SMC})، وتطبيق دالة التشبع (\textenglish{Saturation}) لتخفيف ظاهرة الاهتزاز (\textenglish{Chattering}).
    \item \textbf{المحاكاة الهندسية:} التحقق من استقرار خوارزميات التحكم المطبقة وبناء نموذج محاكاة يحاكي الواقع الفيزيائي عبر ربط بيئة \textenglish{SolidWorks} مع \textenglish{Simscape Multibody} في \textenglish{MATLAB/Simulink}.
    \item \textbf{نظام التفاعل الصوتي:} تصميم طبقة تحكم عالي المستوى (\textenglish{HLC}) تعتمد على حزمة \textenglish{Vosk} لاستخلاص الأوامر الصوتية محلياً دون إنترنت، وترجمتها لحظياً إلى مراجع إحداثية كارتزية توجّه الذراع.
\end{enumerate}

\section{هيكلية الأطروحة}
تم تنظيم هذه الأطروحة في خمسة فصول متسلسلة تعكس المنهجية الهندسية المتبعة في المشروع، وذلك على النحو التالي:
\begin{itemize}
    \item \textbf{الفصل الأول (المقدمة):} يقدم خلفية عامة عن المشروع، دوافعه، أهدافه، والمخطط العام للأطروحة.
    \item \textbf{الفصل الثاني (النمذجة الحركية للذراع الروبوتية):} يتناول توصيف الذراع، تعيين المعلمات الهندسية (\textenglish{DH})، واشتقاق معادلات الحركيات الأمامية والعكسية.
    \item \textbf{الفصل الثالث (النمذجة الديناميكية للذراع الروبوتية):} يعرض حسابات العطالة والجاذبية للحصول على معادلات الحركة، إلى جانب بناء نموذج المحاكاة الميكانيكي للذراع.
    \item \textbf{الفصل الرابع (تصميم أنظمة التحكم للذراع الروبوتية):} يناقش إغلاق حلقات التحكم العالي والمنخفض المستوى، ويستعرض مقارنة تحليلية مفصلة بين خوارزميات التحكم المطبقة للوصول إلى الأداء الديناميكي الأمثل.
    \item \textbf{الفصل الخامس (تصميم نظام التفاعل الصوتي للذراع الروبوتية):} يشرح البنية العتادية والبرمجية لنظام استخلاص الأوامر الصوتية، وآلية ربط الكلمات المفتاحية بالمتجهات الحركية الكارتزية للذراع.
\end{itemize}